\documentclass[10pt,a4paper]{amsart}
\usepackage[margin=19mm]{geometry}
\usepackage{amsmath,amssymb,mathtools}
\usepackage[scr=rsfs,cal=euler]{mathalfa}
\usepackage{xfrac}
\usepackage[unicode,hidelinks,bookmarks=false]{hyperref}
\newcommand{\ie}{i.e.\ }
\newcommand{\cf}{cf.\ }
\newcommand{\resp}{respectively\ }
\newcommand{\Subsection}{\subsection}
\providecommand{\nobreakdash}{\nobreak\hbox{-}}
\setlength{\emergencystretch}{2em}
\multlinegap0pt
\usepackage{amssymb}
\usepackage{mathtools}
\usepackage{xcolor}
\newtheorem{proposition}{Proposition}
\newcommand{\AIP}{\mathsf{AIP}}
\newcommand{\Z}{\mathbb Z}
\newcommand{\bu}{\mathbf{u}}
\newcommand{\minion}[1]{{\mathcal #1}}
\title{Linear equations mod $n$ are pseudo-telepathic}
\author{Lorenzo Ciardo}
\date{Source: arXiv:2609.14632v1 (CC BY 4.0)}
\begin{document}
\maketitle

\noindent\emph{Source and licence.} Linear equations mod $n$ are pseudo-telepathic. Version arXiv:2609.14632v1 (CC BY 4.0), distributed by arXiv under the Creative Commons Attribution 4.0 International licence, http://creativecommons.org/licenses/by/4.0/, as recorded in the arXiv licence metadata for this exact version and shown on the abstract page https://arxiv.org/abs/2609.14632v1. Full licence notice: \texttt{licenses/arxiv-2609.14632-CC-BY-4.0.txt}.\par\smallskip

\noindent\textit{Editorial scope.} Counted unit: the article's proposition prop\_measure (source0.tex lines 328--333). The article's complete original proof of it is delivered with it (source0.tex lines 335--353, one \texttt{proof} environment of 19 lines). No mathematics is written, repaired or re-derived: the statement and the proof are the article's own lines, in the article's own notation, and no retained line is substituted or re-flowed.  No further source-local context is needed: every cross-reference of every retained line resolves inside the two delivered spans.  Nothing was omitted from any retained span, so the package carries no omission record.  No retained line contains a citation, so the wrapper carries no bibliography and no bibliography entry.  No retained line contains a figure environment, an image-inclusion command or drawing code, so no image is delivered and the package declares no asset dependency.  Editorial formatting only, in addition: an AMS \texttt{amsart} preamble that restates the article's own declarations and macros used by the retained lines, namely the environments \texttt{proposition} on separate counters, together with the standard packages those lines need.  The retained environments print in this document as Proposition 1 in order of appearance, whereas the article prints the same items with the numbering of its own full document.  The complete native source of the article as distributed in the arXiv e-print for this exact version is bundled unchanged as \texttt{sources/210338/source0.tex}.
\begin{proposition}
\label{prop_measure}
For every $n\geq2$ and every finite-dimensional Hilbert space $H$, a minion homomorphism $\minion Q_H\to\AIP_n$ exists
if, and only if, there is a $\Z_n$-measure $\mu$ on the subspaces of $H$ satisfying
$\mu(H)=1$.
\end{proposition}

\begin{proof}
Let $\xi:\minion Q_H\to\AIP_n$ be a minion homomorphism, and define $\mu(U)=\xi((U,U^\perp))_1$ for each $U\in L_H$. For any $U\perp V\in L_H$, write $\xi((U,V,(U\oplus V)^\perp))=(a,b,1-a-b)$. Since $\xi$ is a minion homomorphism, we have that $\mu(U)=a$, $\mu(V)=b$, and $\mu(U\oplus V)=a+b$. Hence, $\mu$ is a $\Z_n$-measure on the subspaces of $H$. Moreover, $\xi(H)=1$ and, thus, $\xi((H,\{0\}))=(1,0)$, which means that $\mu(H)=1$.

Conversely, suppose $\mu$ is a $\Z_n$-measure  on the subspaces of $H$ satisfying $\mu(H)=1$. For a tuple $\bu=(U_1,\dots,U_r)\in\minion Q_H^{(r)}$, we let $\xi(\bu)=(\mu(U_1),\dots,\mu(U_r))\in\Z_n^r$. Since the $U_i$'s are mutually orthogonal, the fact that $\mu$ is a measure on the subspaces of $H$ implies that 
    \begin{align*}\sum_{i\in[r]}\mu(U_i)=\mu\left(\bigoplus_{i\in[r]}U_i\right)=\mu(H)=1,
    \end{align*}
    so $\xi(\bu)\in\AIP_n^{(r)}$. To show that $\xi$ is a minion homomorphism, take a map $\pi:[r]\to[s]$, and observe that
    \begin{align*}
        \xi(\bu_{/\pi})
        &=
        \xi\left(\bigoplus_{i\in\pi^{-1}(1)}U_i,\dots,\bigoplus_{i\in\pi^{-1}(s)}U_i\right)
        =
        \left(\mu\left(\bigoplus_{i\in\pi^{-1}(1)}U_i\right),\dots,\mu\left(\bigoplus_{i\in\pi^{-1}(s)}U_i\right)\right)\\
        &=\left(\sum_{i\in\pi^{-1}(1)}\mu(U_i),\dots,\sum_{i\in\pi^{-1}(s)}\mu(U_i)\right)
        =
        \xi(\bu)_{/\pi},
    \end{align*}
    where the third equality uses that subspaces with distinct indices are mutually orthogonal and, thus, $\mu$ behaves additively on them. 
\end{proof}

\end{document}
